% Interim LaTeX file for "Wind Forecasting with Complex ARIMA".
% Replace with the authors' source once uploaded; keep notation consistent with the notes below.
\documentclass[11pt]{article}
\usepackage[margin=2.5cm]{geometry}
\usepackage{amsmath,amssymb,booktabs,graphicx,hyperref}
\usepackage[round]{natbib}
\graphicspath{{../figures/}}

\title{Wind Forecasting with Complex ARIMA}
\author{Sergey Svetunkov \and Pierre Pinson \and Ivan Svetunkov}
\date{Interim draft, \today}

\begin{document}
\maketitle

\begin{abstract}
Abstract text.
\end{abstract}

\section{Introduction}
% Text as in the current PDF draft.

\section{Literature review}

\section{Methodology}
\subsection{Complex representation of wind}
Let $v_t \ge 0$ be the wind speed and $\theta_t$ the wind direction at time $t$, with the
direction mapped to $(-\pi, \pi]$: $\theta_t \leftarrow \theta_t - 2\pi$ whenever $\theta_t > \pi$.
The wind is represented by the complex variable
\begin{equation}
  z_t = v_t e^{i\theta_t} = v_t \cos\theta_t + i\, v_t \sin\theta_t ,
\end{equation}
whose real and imaginary parts are the Cartesian coordinates of the wind vector
(not the zonal and meridional components). In logarithms,
\begin{equation}
  \ln z_t = \ln v_t + i\,\theta_t .
\end{equation}

\subsection{Complex ARIMA}
Following \citet{Svetunkov2024}, cARIMA$(p,d,q)$ is
\begin{equation}
  (1-\phi_1 B-\dots-\phi_p B^p)(1-B)^d y_t = \beta_0 + (1+\vartheta_1 B+\dots+\vartheta_q B^q)\epsilon_t ,
\end{equation}
with complex parameters $\phi_j$, $\vartheta_j$, $\beta_0$ and complex normal $\epsilon_t$.

\subsection{Benchmarks}

\section{Evaluation}
\subsection{Data}
Kaggle Wind Turbine SCADA data: one turbine, 2018, 10-minute resolution.

\subsection{Experimental design}

\subsection{Results}

\section{Conclusions}

\bibliographystyle{apalike}
\bibliography{references}
\end{document}
